\documentclass[11pt]{article}

\usepackage[margin=1in]{geometry}
\usepackage{amsmath}
\usepackage{amssymb}
\usepackage{bm}
\usepackage{graphicx}
\usepackage{array}
\usepackage{booktabs}
\usepackage{xcolor}
\usepackage{mathtools}
\usepackage{enumitem}
\usepackage[numbers,sort&compress]{natbib}
\usepackage{hyperref}
\IfFileExists{cleveref.sty}{\usepackage[nameinlink,noabbrev]{cleveref}}{}

\hypersetup{
  colorlinks=true,
  linkcolor=blue!55!black,
  citecolor=blue!55!black,
  urlcolor=blue!55!black,
  pdftitle={Two-Tracer Redshift-Space Power Spectra in FOLPS},
  pdfauthor={FOLPS cross-power implementation notes}
}

\newcommand{\dd}{\mathrm{d}}
\newcommand{\ii}{\mathrm{i}}
\newcommand{\bk}{\bm{k}}
\newcommand{\bp}{\bm{p}}
\newcommand{\br}{\bm{r}}
\newcommand{\bx}{\bm{x}}
\newcommand{\bn}{\hat{\bm{n}}}
\newcommand{\PL}{P_{\mathrm L}}
\newcommand{\PK}{P_{\mathrm K}}
\newcommand{\PD}{\delta_{\mathrm D}}
\newcommand{\intp}{\int\!\frac{\dd^3\bp}{(2\pi)^3}}
\newcommand{\Afull}{A_{\mathrm{full}}}
\newcommand{\bsst}{b_{s^2}}
\newcolumntype{L}[1]{>{\raggedright\arraybackslash}p{#1}}

\title{\textbf{Two-Tracer Redshift-Space Power Spectra in FOLPS}\\
\large Theory, implementation, and validation notes for equal-time cross spectra}
\author{FOLPS cross-power development notes}
\date{July 25, 2026}

\begin{document}
\maketitle

\begin{abstract}
These notes describe the equal-time two-tracer redshift-space power-spectrum
extension implemented in FOLPS.  The main motivating application is the
cross-correlation of distinct DESI tracers in an overlapping redshift interval,
with LRG3 $\times$ ELG1 as the concrete use case.  The illustrative numerical
figures use LRG-like and ELG-like parameter choices only; they are not DESI
best fits.  The derivation keeps the two tracer endpoints ordered until the
one-loop contractions have been polarized, identifies the exact full biased
$A$-term endpoint factors, and records the implementation and validation checks
used by the current branch.
\end{abstract}

\tableofcontents

\section{Motivation and scope}

Overlapping galaxy samples contain information that is not exhausted by their
two auto spectra.  If tracer $A$ and tracer $B$ occupy a common redshift
interval, the same large-scale matter and velocity fields source both samples,
while their bias expansions, stochastic fields, and pair-level EFT parameters
remain different.  The cross spectrum $P_{AB}$ probes this shared response
directly.  It is also the object that enters the covariance matrix of the
combined two-tracer data vector.

The cross spectrum is not generally reconstructible from $P_{AA}$ and
$P_{BB}$.  The deterministic linear limit has
$P_{AB}^{L}=\sqrt{P_{AA}^{L}P_{BB}^{L}}$, but one-loop bias contractions are
bilinear in the two endpoint bias vectors, EFT counterterms are pair-level
renormalized parameters, and cross stochasticity is not fixed by either auto
spectrum.  A direct one-loop calculation is therefore required.

The first FOLPS implementation is deliberately narrow.  Throughout this note we
assume:
\begin{itemize}[leftmargin=1.6em]
  \item equal time and one common effective redshift for the pair;
  \item common matter and velocity fields, with no tracer velocity bias;
  \item the plane-parallel approximation and a common line of sight $\bn$;
  \item even power-spectrum multipoles only;
  \item the standard EFT redshift-space power spectrum;
  \item optional FolpsD phenomenological cross damping using the prescriptions
        in Sec.~\ref{sec:cross-damping};
  \item no window convolution, covariance, likelihood, or redshift averaging.
\end{itemize}

This implementation branch is based directly on the repository \texttt{main}
branch.  It is not based on the separate \texttt{adematti-damping}
development branch, and it does not claim feature parity with that branch.
Damping support is limited to functionality present and tested on this
\texttt{main}-based branch.  Reconciling this work with other
damping-development branches is outside the present scope and is not a
scientific dependency of the cross-spectrum theory described here.

Under these assumptions the existing FOLPS FFTLog loop tables are sufficient.
The extension reuses the matter, velocity, bias, $A$, $D$, $G$, no-wiggle, and
IR-resummation rows, replacing each single-tracer polynomial by the correctly
polarized two-tracer polynomial.

\section{Redshift-space definition}

Let $\bn$ be the fixed line-of-sight direction and define
$\mu=\hat{\bk}\cdot\bn$.  The redshift-space cross spectrum is defined by
\begin{equation}
  \left\langle
    \delta_A^s(\bk)\delta_B^s(\bk')
  \right\rangle
  =
  (2\pi)^3\PD(\bk+\bk')P^s_{AB}(k,\mu).
\end{equation}
The exact equal-time plane-parallel generating function is
\begin{align}
  (2\pi)^3\PD(\bk)+P^s_{AB}(\bk)
  ={}&
  \int \dd^3\br\,e^{-\ii\bk\cdot\br}
  \left\langle
  [1+\delta_A(\bx_1)]
  [1+\delta_B(\bx_2)]
  e^{-\ii k_\parallel\Delta v_\parallel/(aH)}
  \right\rangle ,
  \label{eq:generating}
\end{align}
where $\br=\bx_2-\bx_1$,
$k_\parallel=\bk\cdot\bn$, and
\begin{equation}
  \Delta v_\parallel=
  \bn\cdot[\bm{v}(\bx_2)-\bm{v}(\bx_1)] .
\end{equation}
The fields $\delta_A$ and $\delta_B$ denote the real-space density contrasts of
the two tracers.  Their labels are endpoint labels in
\eqref{eq:generating}; they must not be identified until after the
contractions have been evaluated.

At equal time, with statistical homogeneity, isotropy, parity invariance, and a
common plane-parallel line of sight, the spectrum is real and symmetric under
endpoint exchange:
\begin{equation}
  P^s_{AB}(k,\mu)=P^s_{BA}(k,\mu).
  \label{eq:exchange}
\end{equation}
The same assumptions give $P^s_{AB}(k,\mu)=P^s_{AB}(k,-\mu)$, so the present
implementation projects only even multipoles,
\begin{equation}
  P^{AB}_\ell(k)=
  \frac{2\ell+1}{2}\int_{-1}^{1}\dd\mu\,
  \mathcal{L}_\ell(\mu)P^s_{AB}(k,\mu),
  \qquad \ell=0,2,4.
\end{equation}
Odd multipoles can arise once wide-angle, relativistic, unequal-time, or
selection-asymmetry effects are retained; those effects are outside this first
FOLPS cross-power path.

\section{Bias conventions}

The public prior-document input basis is
\begin{equation}
  b_1,\qquad b_2,\qquad b_{K^2},\qquad b_{\rm td}.
\end{equation}
The current FOLPS power-spectrum API maps these four entries to the canonical
internal basis
\begin{equation}
  b_{s^2}=2b_{K^2},
  \qquad
  b_{3\mathrm{nl}}
  =
  -\frac{32}{21}
  \left(
    b_{K^2}+\frac{2}{5}b_{\rm td}
  \right).
  \label{eq:priordoc-map}
\end{equation}
In canonical FOLPS parameters the density field of tracer $X$ is represented
schematically as
\begin{equation}
  \delta_X
  =
  b_1^X\delta
  +\frac{b_2^X}{2}[\delta^2]
  +\frac{\bsst^X}{2}[s^2]_{\rm code}
  +b_{3\mathrm{nl}}^X O_{3\mathrm{nl}}
  +\epsilon_X+\cdots .
\end{equation}
The explicit factor $1/2$ in the $b_2$ source is the standard local-quadratic
convention.  The current FOLPS full-$A$ kernels also insert canonical
$b_{s^2}$ with an explicit factor $1/2$ in the same second-order source block.
Consequently, the full biased $A$ cross term has one factor $1/2$ from the
second-order source convention and a second factor $1/2$ from polarizing an
endpoint-summed auto row.  This is the origin of the factors $1/4$ in
Sec.~\ref{sec:Aterm}.

\section{One-loop decomposition}

FOLPS organizes the equal-time cross spectrum as
\begin{align}
  P_{AB}^{s}(k,\mu)
  ={}&
  P_{\delta_A\delta_B}
  +f_0\mu^2
  \left(
    P_{\delta_A\theta}
    +
    P_{\delta_B\theta}
  \right)
  +f_0^2\mu^4P_{\theta\theta}
  \nonumber\\
  &+A_{AB}
  +D_{AB}
  +G_{AB}
  +P_{\mathrm{ctr},AB}
  +P_{\mathrm{stoch},AB}.
  \label{eq:folps-decomp}
\end{align}
Here $\theta$ is the common velocity-divergence field in the FOLPS
normalization, $f_0$ is the large-scale growth rate, and $f(k)$ is the
scale-dependent growth rate used by the FOLPS tables.  At linear order,
\begin{equation}
  P_{\mathrm K,AB}
  =
  \left(b_1^A+f(k)\mu^2\right)
  \left(b_1^B+f(k)\mu^2\right)
  \PL(k).
  \label{eq:kaiser}
\end{equation}
The terms in \eqref{eq:folps-decomp} beyond the cross-Kaiser spectrum are the
one-loop density, velocity, redshift-space mode-coupling, EFT, and stochastic
pieces.

\section{Density and velocity loops}

The biased density-density loop polynomial is obtained by bilinearizing the
single-tracer FOLPS auto-spectrum polynomial:
\begin{align}
P_{\delta_A\delta_B}^{\rm loop}
={}&b_1^A b_1^B P_{\delta\delta}^{\rm loop}
+(b_1^A b_2^B+b_2^A b_1^B)P_{b_1b_2}
\nonumber\\
&+(b_1^A \bsst^B+\bsst^A b_1^B)P_{b_1s^2}
+b_2^A b_2^B P_{b_2b_2}
\nonumber\\
&+(b_2^A \bsst^B+\bsst^A b_2^B)P_{b_2s^2}
+\bsst^A \bsst^B P_{s^2s^2}
\nonumber\\
&+(b_1^A b_{3\mathrm{nl}}^B+b_{3\mathrm{nl}}^A b_1^B)
\sigma_3^2\PL .
\label{eq:density-loop}
\end{align}
For either tracer $X\in\{A,B\}$,
\begin{equation}
P_{\delta_X\theta}^{\rm loop}
=
b_1^X P_{\delta\theta}^{\rm loop}
+b_2^X P_{b_2\theta}
+\bsst^X P_{s^2\theta}
+b_{3\mathrm{nl}}^X\frac{f(k)}{f_0}\sigma_3^2\PL .
\label{eq:density-velocity-loop}
\end{equation}
The velocity-velocity loop $P_{\theta\theta}^{\rm loop}$ is common to both
tracers.  Equations \eqref{eq:density-loop} and
\eqref{eq:density-velocity-loop} use only loop rows that already exist in the
FOLPS tables; no new FFTLog matrices are introduced.

\section{The full biased A term}
\label{sec:Aterm}

The $A$ term is the part of the redshift-space expansion with one pairwise
velocity insertion and one ordered bispectrum.  Define
\begin{equation}
  \mathcal{S}_X(\bm{q})
  =
  \delta_X(\bm{q})+f_0\mu_q^2\theta(\bm{q})
\end{equation}
and the ordered bispectrum
\begin{equation}
  (2\pi)^3\PD(\bk_1+\bk_2+\bk_3)
  B_{\sigma;AB}(\bk_1,\bk_2,\bk_3)
  =
  \left\langle
    \theta(\bk_1)\mathcal{S}_A(\bk_2)\mathcal{S}_B(\bk_3)
  \right\rangle .
\end{equation}
The TNS/FOLPS $A$ term can be written as
\begin{equation}
  A_{AB}(k,\mu)=
  k\mu f_0\intp
  \frac{p_\parallel}{p^2}
  \left[
  B_{\sigma;AB}(\bp,\bk-\bp,-\bk)
  -
  B_{\sigma;AB}(\bp,\bk,-\bk-\bp)
  \right].
\end{equation}

\subsection{Endpoint exchange}

The Stage 1B endpoint audit resolves the central matrix question.  A P22-type
FFTLog scalar has the form
\begin{equation}
  S[x,M,y]=K^3\sum_{ij}x_iM_{ij}y_j .
\end{equation}
Endpoint exchange is the loop-variable transformation
$\bp\leftrightarrow\bk-\bp$ together with exchange of the two endpoint spectra.
Therefore the physical exchange is
\begin{equation}
  x^{\mathsf T}My
  \longrightarrow
  y^{\mathsf T}M^{\mathsf T}x
  =
  x^{\mathsf T}My .
  \label{eq:endpoint-exchange}
\end{equation}
The exchanged endpoint is not represented by $x^{\mathsf T}M^{\mathsf T}y$.
That object transposes the kernel while leaving the endpoint vectors attached
to the wrong endpoints, and is generally different when the vectors differ.
Equation \eqref{eq:endpoint-exchange} is algebraic, not Hermitian; no complex
conjugation is involved.

\subsection{Reduced A term}

With \texttt{A\_full=False}, the exact cross polynomial is
\begin{align}
A_{AB}^{\rm reduced}
={}&b_1^A b_1^B f_0\mu^2 I_{1udd,1}
\nonumber\\
&+\frac{b_1^A+b_1^B}{2}f_0^2
\left(
\mu^2I_{2uud,1}
+\mu^4I_{2uud,2}
\right)
\nonumber\\
&+f_0^3
\left(
\mu^4I_{3uuu,2}
+\mu^6I_{3uuu,3}
\right).
\label{eq:A-reduced}
\end{align}
The mixed $I_{2uud}$ rows are endpoint-summed auto rows, so they are polarized
by $(b_1^A+b_1^B)/2$.  The $I_{1udd,1}$ and pure velocity rows are already
symmetric under endpoint exchange.

\subsection{Additional full-A rows}

With \texttt{A\_full=True}, the biased second-order rows add
\begin{align}
\Delta A_{AB}^{b_2}
={}&
\frac{b_2^A b_1^B+b_1^A b_2^B}{4}
f_0\mu^2I_{1udd,1}^{b_2}
\nonumber\\
&+
\frac{b_2^A+b_2^B}{4}f_0^2
\left(
\mu^2I_{2uud,1}^{b_2}
+\mu^4I_{2uud,2}^{b_2}
\right),
\label{eq:A-b2}
\end{align}
and
\begin{align}
\Delta A_{AB}^{s^2}
={}&
\frac{\bsst^A b_1^B+b_1^A \bsst^B}{4}
f_0\mu^2I_{1udd,1}^{s^2}
\nonumber\\
&+
\frac{\bsst^A+\bsst^B}{4}f_0^2
\left(
\mu^2I_{2uud,1}^{s^2}
+\mu^4I_{2uud,2}^{s^2}
\right).
\label{eq:A-bs2}
\end{align}
Each factor $1/4$ is the product of two independent factors: $1/2$ from the
second-order source convention and $1/2$ from polarizing an endpoint-summed
auto row.  In the auto limit $A=B$, \eqref{eq:A-b2} and \eqref{eq:A-bs2}
reduce exactly to the canonical FOLPS auto coefficients.

\begin{table}[p]
\centering
\footnotesize
\setlength{\tabcolsep}{3pt}
\caption{Current FOLPS matrix and loop-table mapping for the $A$ rows.  The
suffix ``\_b'' on vectors denotes the auxiliary FFTLog bias exponent used by
FOLPS, not tracer $B$.}
\label{tab:A-mapping}
\begin{tabular}{@{}L{0.18\textwidth}L{0.24\textwidth}L{0.13\textwidth}L{0.13\textwidth}L{0.19\textwidth}@{}}
\toprule
Function & Matrix or vector pieces & Left vector & Right vector & Loop-table row \\
\midrule
$I_{1udd,1}$ & $M^{tAfp}_{11}$ plus $M^{afk}_{11},M^{afp}_{11}$ & \texttt{vecf} & \texttt{vec} & \texttt{I1udd\_1} \\
$I_{2uud,1}$ & $M^{tAfkmpfp}_{12}$ plus $M^{afkfp}_{12},M^{afpfp}_{12}$ & \texttt{vecf} & \texttt{vecf} & \texttt{I2uud\_1} \\
$I_{2uud,2}$ & $M^{tAfkmpfp}_{22}$, $M^{tAfpfp}_{22}$ and P13 pieces & \texttt{vecf}, \texttt{vecff} & \texttt{vecf}, \texttt{vec} & \texttt{I2uud\_2} \\
$I_{3uuu,2}$ & $M^{tAfkmpfpfp}_{23}$ plus P13 inherited pieces & \texttt{vecff} & \texttt{vecf} & \texttt{I3uuu\_2} \\
$I_{3uuu,3}$ & $M^{tAfkmpfpfp}_{33}$ plus $M^{afkfkfp}_{33},M^{afkfpfp}_{33}$ & \texttt{vecff} & \texttt{vecf} & \texttt{I3uuu\_3} \\
$I_{1udd,1}^{b_2}$ & $M^{tAfp}_{11,b2}$ & \texttt{vecf\_b} & \texttt{vec\_b} & \texttt{I1udd\_1\_b2} \\
$I_{2uud,1}^{b_2}$ & $M^{tAfkmpfp}_{12,b2}$ & \texttt{vecf\_b} & \texttt{vecf\_b} & \texttt{I2uud\_1\_b2} \\
$I_{2uud,2}^{b_2}$ & $M^{tAfkmpfp}_{22,b2}$ & \texttt{vecf} & \texttt{vecf} & \texttt{I2uud\_2\_b2} \\
$I_{1udd,1}^{s^2}$ & $M^{tAfp}_{11,bs2}$ & \texttt{vecf\_b} & \texttt{vec\_b} & \texttt{I1udd\_1\_bs2} \\
$I_{2uud,1}^{s^2}$ & $M^{tAfkmpfp}_{12,bs2}$ & \texttt{vecf\_b} & \texttt{vecf\_b} & \texttt{I2uud\_1\_bs2} \\
$I_{2uud,2}^{s^2}$ & $M^{tAfkmpfp}_{22,bs2}$ & \texttt{vecf} & \texttt{vecf} & \texttt{I2uud\_2\_bs2} \\
\bottomrule
\end{tabular}
\end{table}

\section{D and G terms}

The cross $D$ term uses the same row basis as the auto spectrum.  Define
\begin{align}
D_2 &=
\mu^2 I^D_{2uudd,1}
+\mu^4 I^D_{2uudd,2},
\nonumber\\
D_3 &=
\mu^2 I^B_{3uuud,1}
+\mu^4 I^D_{3uuud,2}
+\mu^6 I^D_{3uuud,3},
\nonumber\\
D_4 &=
\mu^2 I^B_{4uuuu,1}
+\mu^4 I^D_{4uuuu,2}
+\mu^6 I^D_{4uuuu,3}
+\mu^8 I^D_{4uuuu,4}.
\end{align}
Then
\begin{equation}
D_{AB}
=
b_1^A b_1^B f_0^2D_2
+
\frac{b_1^A+b_1^B}{2}f_0^3D_3
+
f_0^4D_4 .
\label{eq:D-cross}
\end{equation}
The $D_3$ coefficient is polarized because the single density insertion can be
assigned to either endpoint.

The $G$ term is proportional to the cross-Kaiser polynomial in the current
FOLPS convention:
\begin{equation}
G_{AB}
=
-(k\mu f_0)^2\sigma_{2w}
\left[
b_1^A b_1^B \PL
+
(b_1^A+b_1^B)f_0\mu^2P_{\delta\theta}^{L}
+
f_0^2\mu^4P_{\theta\theta}^{L}
\right],
\label{eq:G-cross}
\end{equation}
where
$P_{\delta\theta}^{L}=(f(k)/f_0)\PL$ and
$P_{\theta\theta}^{L}=(f(k)/f_0)^2\PL$.

\section{EFT and stochastic terms}

The first cross implementation uses explicit pair-level nuisance parameters:
\begin{align}
P_{\mathrm{ctr},AB}
&=
(\alpha_0^{AB}
+\alpha_2^{AB}\mu^2
+\alpha_4^{AB}\mu^4)
k^2\PL,
\label{eq:ctr}
\\
P_{\mathrm{ctr,NLO},AB}
&=
\widetilde c^{AB}
(k\mu f_0)^4
\sigma_{2w}^2
P_{\mathrm K,AB},
\label{eq:ctr-nlo}
\\
P_{\mathrm{stoch},AB}
&=
P_{\mathrm{shot}}^{AB}
\left[
\alpha_{\mathrm{shot},0}^{AB}
+
\alpha_{\mathrm{shot},2}^{AB}(k\mu)^2
\right].
\label{eq:stoch}
\end{align}
These parameters are not inferred from the two auto-spectrum nuisance vectors.
The public cross API requires them explicitly.  The stochastic term is not
identified with a Poisson-overlap prediction; $P_{\mathrm{shot}}^{AB}$ is a
pair-level stochastic normalization in the current parameterization.  The
implementation does not currently impose a positivity-preserving covariance
parametrization for the joint $AA$, $AB$, $BB$ model.

\subsection{Phenomenological cross damping}
\label{sec:cross-damping}

For \texttt{model="FOLPSD"} the cross-power API supports the same production
power-spectrum damping names as the auto path, \texttt{"exp"}, \texttt{"lor"},
and \texttt{"vdg"}.  The public keyword
\texttt{cross\_damping\_mode} chooses how the cross damping factor is sourced.
The default is \texttt{"single"},
\begin{equation}
W_{AB}^{\mathrm{single}}(k,\mu)
=
W(k,\mu;X_{\rm FoG}^{AB}),
\end{equation}
where $X_{\rm FoG}^{AB}$ is the final entry of the explicit pair-level
\texttt{cross\_nuisance} vector.  This introduces one damping nuisance
parameter for the measured cross spectrum.

The alternative mode is \texttt{"geometric"},
\begin{equation}
W_{AB}^{\mathrm{geometric}}(k,\mu)
=
\sqrt{
W(k,\mu;X_{\rm FoG}^{A})
W(k,\mu;X_{\rm FoG}^{B})
}.
\end{equation}
Here $X_{\rm FoG}^{A}$ and $X_{\rm FoG}^{B}$ are read from the two tracer
parameter arrays after each has been converted through its own bias convention.
The pair-level $X_{\rm FoG}^{AB}$ remains present in
\texttt{cross\_nuisance} for API compatibility, but it is ignored in this mode.
The geometric prescription preserves exchange symmetry because swapping the
tracers only interchanges the two factors under the square root.  When
$X_{\rm FoG}^{A}=X_{\rm FoG}^{B}$, it reduces to the corresponding auto damping
factor.

The damping placement is intentionally identical to the current auto FolpsD
implementation.  In the production non-marginalized code, $W$ multiplies the
one-loop pair contraction assembled inside \texttt{PloopSPTs\_cross}.  The
IR-resummed linear Kaiser term, the standard counterterms, the NLO
$\widetilde c$ counterterm, and the stochastic term are not multiplied by this
phenomenological damping factor under the current \texttt{Winfty\_all=False}
convention.  Counterterms and stochastic terms retain the pair-level
parameterization in Eqs.~\eqref{eq:ctr}--\eqref{eq:stoch}.

The two damping prescriptions are alternative nuisance models.  Neither is
claimed to be uniquely selected by first principles; mocks and data should be
used to decide whether the extra tracer-level structure in the geometric mode
is useful.

\begin{figure}[p]
  \centering
  \includegraphics[width=\textwidth]{figures/cross_damping_modes.pdf}
  \caption{Illustrative FolpsD cross damping comparison for VDG damping.  The
  single mode uses one pair-level $X_{\rm FoG}^{AB}$, while the geometric mode
  uses the two tracer-level damping factors.  The lower panels show the
  fractional difference relative to the single prescription.}
  \label{fig:cross-damping-modes}
\end{figure}

\section{Why the geometric mean is not enough}
\label{sec:geom}

In the deterministic linear limit,
\begin{equation}
  P_{AB}^{L}
  =
  \sqrt{P_{AA}^{L}P_{BB}^{L}},
\end{equation}
because all spectra are products of the same linear matter power spectrum and
linear endpoint responses.  Beyond that limit, write schematically
\begin{equation}
  P_{XX}=P_{XX}^{L}+P_{XX}^{1\text{-loop}} .
\end{equation}
Then
\begin{equation}
\sqrt{P_{AA}P_{BB}}
\simeq
\sqrt{P_{AA}^{L}P_{BB}^{L}}
\left[
1+\frac12
\frac{P_{AA}^{1\text{-loop}}}{P_{AA}^{L}}
+\frac12
\frac{P_{BB}^{1\text{-loop}}}{P_{BB}^{L}}
\right],
\label{eq:geom-expansion}
\end{equation}
which averages fractional auto corrections.  The true cross loop instead
contains bilinear monomials such as $b_2^A b_1^B+b_1^A b_2^B$ and
$b_2^A\bsst^B+\bsst^A b_2^B$.  Cross EFT and stochastic terms are independent
pair-level renormalized parameters.  Therefore the geometric mean is at most a
diagnostic approximation.

\begin{figure}[p]
  \centering
  \includegraphics[width=0.86\textwidth]{figures/cross_vs_geometric_mean_pkmu.pdf}
  \caption{Comparison of the implemented $P^{AB}(k,\mu)$ with the diagnostic
  geometric mean $\sqrt{P^{AA}(k,\mu)P^{BB}(k,\mu)}$ for three angular slices.
  The lower panel displays $P^{AB}/\sqrt{P^{AA}P^{BB}}-1$ only where the
  geometric mean is real and numerically stable.}
  \label{fig:geom-pkmu}
\end{figure}

\begin{figure}[p]
  \centering
  \includegraphics[width=0.86\textwidth]{figures/cross_geometric_mean_bias_dependence.pdf}
  \caption{Controlled sector switch for the geometric-mean diagnostic at
  $\mu=0.5$.  The deterministic linear synthetic table gives $r-1$ consistent
  with zero, while matter/RSD loops, nonlinear bias terms, and explicit
  pair-level nuisance parameters generate departures from the geometric mean.}
  \label{fig:geom-bias}
\end{figure}

\section{IR resummation and AP mapping}

The cross implementation keeps the existing FOLPS IR-resummation structure.
The displacement damping is shared because the two tracers live in the same
matter and velocity fields.  Only the tracer-dependent polynomials are replaced
by their cross forms.  In code notation,
\begin{align}
\Sigma_t^2(k,\mu)
&=
\left[1+f_0\mu^2(2+f_0)\right]\Sigma^2
+(f_0\mu)^2(\mu^2-1)\Delta\Sigma^2 ,
\nonumber\\
E(k,\mu)&=\exp[-k^2\Sigma_t^2(k,\mu)] ,
\end{align}
and the returned redshift-space spectrum is assembled as
\begin{align}
P_{AB}^{s,\mathrm{IR}}
={}&
P_{\mathrm K,AB}
\left[
P_{\mathrm{nw}}
+E(P_{\mathrm L}-P_{\mathrm{nw}})(1+k^2\Sigma_t^2)
\right]
\nonumber\\
&+
E\,P_{AB}^{\rm 1loop}[\mathrm{wiggle}]
+
(1-E)\,P_{AB}^{\rm 1loop}[\mathrm{nowiggle}] .
\end{align}

The Alcock-Paczynski mapping is also unchanged.  With
$F=q_\parallel/q_\perp$,
\begin{align}
k_{\rm true}
&=
\frac{k_{\rm obs}}{q_\perp}
\left[
1+\mu_{\rm obs}^2\left(F^{-2}-1\right)
\right]^{1/2},
\nonumber\\
\mu_{\rm true}
&=
\frac{\mu_{\rm obs}}{F}
\left[
1+\mu_{\rm obs}^2\left(F^{-2}-1\right)
\right]^{-1/2},
\end{align}
and the multipoles are multiplied by the Jacobian
\begin{equation}
  (q_\parallel q_\perp^2)^{-1}.
\end{equation}

\begin{figure}[p]
  \centering
  \includegraphics[width=\textwidth]{figures/cross_ir_resummation.pdf}
  \caption{Numerical effect of the existing FOLPS IR resummation on the cross
  multipoles.  The upper row compares the resummed and non-resummed
  calculations; the lower row shows their absolute difference.}
  \label{fig:ir}
\end{figure}

\begin{figure}[p]
  \centering
  \includegraphics[width=\textwidth]{figures/cross_ap_remapping.pdf}
  \caption{Illustrative Alcock-Paczynski remapping for the cross multipoles.
  The distorted case uses $q_\parallel=1.02$ and $q_\perp=0.98$ as a numerical
  demonstration, not as a DESI fiducial choice.}
  \label{fig:ap}
\end{figure}

\section{Implementation architecture}

The implementation remains in \texttt{folps/folps.py}; no \texttt{folpsX.py}
module is introduced.  The non-marginalized auto and cross spectra both pass
through the same private pair-level helper.  Auto mode is obtained by setting
$A=B$ and using the auto nuisance tuple.  Cross mode is activated by passing
\texttt{pars\_b} and an explicit \texttt{cross\_nuisance} tuple.

The implementation uses small \texttt{NamedTuple} containers for the canonical
bias and nuisance subsets.  The public multipole method accepts
\texttt{bias\_scheme} for tracer $A$ and \texttt{bias\_scheme\_b} for tracer
$B$; the prior-document convention is converted by \eqref{eq:priordoc-map}
before entering the pair contraction.  Direct \texttt{get\_rsd\_pkmu} and
\texttt{get\_eft\_pkmu} calls use canonical FOLPS parameter ordering.

Cross spectra are supported for \texttt{model="EFT"} and
\texttt{model="FOLPSD"}.  With \texttt{model="FOLPSD"}, the public
\texttt{cross\_damping\_mode} keyword accepts \texttt{"single"} and
\texttt{"geometric"}, with \texttt{"single"} as the default.  The table and
\texttt{table\_now} schemas are reused directly,
including the six \texttt{A\_full=True} rows, the no-wiggle rows, and the IR
scalar tail.  No matrix file is added or modified.  The NumPy and JAX backends
share the same source formulas through the module-selected backend namespace;
JAX device selection is environment dependent and is not restricted to CPU by
the cross-power code.

The FolpsD cross damping path supports only the main-branch damping names
\texttt{"exp"}, \texttt{"lor"}, and \texttt{"vdg"}.  No additional damping
names, fallback aliases, or flags from \texttt{adematti-damping} are part of
this branch.

\begin{verbatim}
pars_a = [...]
pars_b = [...]
cross_nuisance_ab = [
    alpha0_ab, alpha2_ab, alpha4_ab, ctilde_ab,
    alphashot0_ab, alphashot2_ab, PshotP_ab, X_FoG_ab,
]

multipoles = RSDMultipolesPowerSpectrumCalculator(model="FOLPSD")
p0_ab, p2_ab, p4_ab = multipoles.get_rsd_pkell(
    kobs, 1.0, 1.0, pars_a, table, table_now,
    bias_scheme="priordoc",
    bias_scheme_b="priordoc",
    pars_b=pars_b,
    cross_nuisance=cross_nuisance_ab,
    damping="vdg",
    cross_damping_mode="single",
    IR_resummation=True,
    ells=(0, 2, 4),
)
\end{verbatim}

\section{Validation}

The independent tests in \texttt{folps/test\_cross\_power\_spectrum.py}
exercise both algebraic identities and backend behavior:
\begin{itemize}[leftmargin=1.6em]
  \item linear Kaiser recovery with synthetic tables;
  \item one-row synthetic loop checks for every density, density-velocity,
        velocity-velocity, reduced-$A$, full-$A$, $D$, and $G$ row;
  \item all six full-$A$ rows with asymmetric tracer biases;
  \item exchange symmetry $P_{AB}=P_{BA}$;
  \item auto-limit recovery for \texttt{A\_full=False} and \texttt{A\_full=True};
  \item prior-document/canonical equivalence;
  \item explicit missing-\texttt{cross\_nuisance} rejection;
  \item IR on/off checks;
  \item FolpsD cross damping checks for \texttt{"exp"}, \texttt{"lor"}, and
        \texttt{"vdg"} in both \texttt{"single"} and \texttt{"geometric"}
        modes;
  \item AP-remapped multipole checks against direct quadrature;
  \item NumPy/JAX agreement;
  \item unchanged auto-spectrum regression against existing outputs.
\end{itemize}
For the plotted run, the exchange, auto-limit, and prior-document/canonical
residuals are exactly zero.  In Fig.~\ref{fig:residuals}, the $10^{-16}$ level
is only a visualization floor for logarithmic plotting of those zero residuals.

\begin{figure}[p]
  \centering
  \includegraphics[width=\textwidth]{figures/cross_implementation_residuals.pdf}
  \caption{Implementation residuals for exchange symmetry, auto-limit recovery,
  and prior-document versus canonical input conversion.  The logarithmic floor
  is a display device only and has no physical significance.}
  \label{fig:residuals}
\end{figure}

\section{Numerical illustrations}

The generated figure set uses the public input linear spectrum and the existing
\texttt{A\_full=True}, \texttt{use\_TNS\_model=False} matrix cache.  The common
settings are \texttt{model="EFT"}, \texttt{damping=None}, IR resummation on,
even multipoles $(0,2,4)$, and prior-document bias inputs.  All parameter
changes in this section are illustrative visibility choices, not priors.  The
cross damping comparison in Fig.~\ref{fig:cross-damping-modes} explicitly
switches to \texttt{model="FOLPSD"} and \texttt{damping="vdg"}.

\begin{figure}[p]
  \centering
  \includegraphics[width=\textwidth]{figures/cross_power_multipoles.pdf}
  \caption{Illustrative auto and cross multipoles for the LRG-like tracer $A$,
  the ELG-like tracer $B$, and their equal-time cross spectrum.  The cross
  multipole is a projected EFT model output and is not constrained to lie
  between the two auto multipoles at every $k$ and $\ell$.}
  \label{fig:multipoles}
\end{figure}

\begin{figure}[p]
  \centering
  \includegraphics[width=0.86\textwidth]{figures/cross_nonlinear_bias_response.pdf}
  \caption{Fractional response of the cross monopole when one nonlinear
  prior-document bias parameter is varied for tracer $A$ or tracer $B$.  The
  independent responses illustrate that the cross loop polynomial is genuinely
  bilinear in the two tracer bias vectors.}
  \label{fig:bias-response}
\end{figure}

\begin{figure}[p]
  \centering
  \includegraphics[width=\textwidth]{figures/cross_nuisance_response.pdf}
  \caption{Response of the cross multipoles to explicit pair-level EFT and
  stochastic parameters.  The rows isolate the standard EFT counterterms, the
  NLO $\widetilde c^{AB}$ counterterm, and the stochastic terms with
  independent y-axis ranges so the smaller responses remain visible.}
  \label{fig:nuisance-response}
\end{figure}

\section{Limitations and future work}

The current implementation is the first equal-time two-tracer power-spectrum
path.  The main limitations and natural extensions are:
\begin{itemize}[leftmargin=1.6em]
  \item effective-redshift modeling and redshift averaging across the true
        overlap window;
  \item realistic LRG3 $\times$ ELG1 parameter choices and data-vector
        definitions;
  \item cross stochasticity models that enforce positivity of the joint
        $2\times2$ tracer covariance;
  \item field-level counterterm parametrizations that relate auto and cross EFT
        coefficients;
  \item survey window convolution and integral-constraint effects;
  \item covariance and likelihood integration, including desilike interfaces;
  \item tracer velocity bias;
  \item odd multipoles, wide-angle terms, and relativistic corrections;
  \item mock- and data-driven comparisons of the available FolpsD cross
        damping nuisance prescriptions;
  \item bispectrum cross-correlations and joint power-spectrum/bispectrum
        two-tracer analyses.
\end{itemize}

\section{Reference metadata}

Reference metadata were checked against the local notes and arXiv records for
the four required entries.  The original TNS paper is
Ref.~\cite{Taruya2010}; the FOLPS and FFTLog massive-neutrino references are
Refs.~\cite{Aviles2021,Noriega2022}; and the FolpsD reference is
Ref.~\cite{Bansal2026}.

\bibliographystyle{unsrtnat}
\bibliography{source/references}

\end{document}
